\section{Hypothesis}
\label{sec_intro_hypothesis}

The investigation is organised around a single composite hypothesis.

\begin{quote}
    \emph{Multi-person pose grouping is fundamentally an embedding-quality problem rather than a grouping-algorithm problem.
    A skeleton-aware graph attention network trained on synthetic data and fine-tuned 
    on real images can produce per-keypoint embeddings on which a simple unconstrained clustering algorithm
    recovers person identities at quality competitive with an end-to-end joint-trained alternative,
    while supporting autonomous person-count estimation that requires no oracle at inference 
    and a modular composition with any keypoint detector.}
\end{quote}

This hypothesis has three distinct parts: a methodological, an architectural and an operational claim.

\begin{itemize}
    \item \textbf{Methodological claim:} The binding constraint on bottom-up pose grouping is the embeddings rather
    than the algorithm that operates on them. 
    \item \textbf{Architectural claim:} Skeleton-aware graph attention network (GAT) embeddings
    (trained on synthetic data then fine-tuned) are the structural realisation of the methodological claim.
    \item \textbf{Operational claim:} The same embeddings would allow the grouping to be deployed 
    using arbitrary detectors and without an oracle at inference.
\end{itemize}

The three clauses are testable on their own and each decomposes into the specific research questions in the next section.
Chapter~\ref{chap_6} discusses each clause and states which parts of the hypothesis the experimental
evidence of Chapter~\ref{chap_4} supports, and with what caveats.
The research questions are then the instruments by which each clause is interrogated.
Success or failure on any single question neither confirms nor refutes the hypothesis on its own, but
together, the seven questions span its testable surface.